Das \emph{zweite Problem} möchte ich skizzieren im Fall der Darstellung von hyperkomplexen Systemen, allgemeiner von Ringen, die dem \guillemotleft Doppelkettensatz\guillemotright{} genügen. Das Resultat ist, dass jede irreduzible (einfache) Modulklasse Idealklasse ist, dass also die Kompositionsreihen aus (einseitigen) Idealen durch ihre Faktorgruppen schon alle irreduzibeln Darstellungen liefern. Und zwar genügen genauer die Kompositionsreihen im Restklassenring nach dem Radikal, welch letzterer ein \guillemotleft vollständig reduzibler\guillemotright{} Ring wird. Das Problem ist damit zurückgeführt auf die Strukturuntersuchung der vollständig reduzibeln Ringe; und damit erweist es sich als dem Problem entsprechend, die Darstellungen in den Automorphismenkörpern als Ausgangspunkt zu nehmen. Darunter ist das Folgende zu verstehen: Alle einfachen (Rechts) Ideale einer zweiseitig einfachen Komponente eines solchen Ringes gehören derselben Klasse an, besitzen also denselben Automorphismenring, der wegen der Einfachheit der Ideale zum Körper wird, zugleich Automorphismenring der einfachen Linksideale ist. Der einfache Ring selbst wird dem System aller Matrizen im Automorphismenkörper isomorph; und aus dieser Darstellung entspringen alle Darstellungen inbezug auf einen Unterkörper, also insbesondere inbezug auf den Koeffizientenbereich des hyperkomplexen Systems, indem man die Elemente des Automorphismenkörpers selbst durch zugeordnete Matrizen ersetzt. So wird hier der im Fall der hyperkomplexen Systeme bekannte Wedderburnsche Struktursatz durch den der allgemeinen Gruppentheorie entstammenden Begriff des Automorphismenkörpers neu gedeutet, und zugleich als Quelle der Darstellungstheorie der hyperkomplexen Systeme erkannt. Es entstehen die neuen Fragen der Galoisschen Theorie nichtkommutativer Körper, die eng mit der Frage der \guillemotleft Zerfällungskörper\guillemotright{} zusammenhängen. Zugleich ergibt sich ein Ausblick auf Struktursätze allgemeiner, nicht vollständig reduzibler Ringe vermöge der Automorphismenringe der unzerlegbaren Ideale. Die Gruppentheorie erweist sich, in ihrer Ausdehnung auf allgemeinste Gruppen mit Operatoren, auch hier wieder als ordnendes Prinzip.
